\documentclass[11pt]{article}

\usepackage[margin=1.1in]{geometry}
\usepackage{amsmath,amssymb,amsthm,mathtools}
\usepackage{enumitem}
\usepackage{hyperref}
\usepackage{microtype}
\usepackage{mathrsfs}

\hypersetup{
  colorlinks=true,
  linkcolor=blue,
  citecolor=blue,
  urlcolor=blue,
  hypertexnames=false,
  pdftitle={Poisson--Dirichlet Universality for Luce Permutations},
  pdfauthor={Christopher D. Long}
}

\newtheorem{theorem}{Theorem}[section]
\newtheorem{proposition}[theorem]{Proposition}
\newtheorem{lemma}[theorem]{Lemma}
\newtheorem{corollary}[theorem]{Corollary}
\theoremstyle{remark}
\newtheorem{remark}[theorem]{Remark}

\newcommand{\Pp}{\mathbb P}
\newcommand{\Ee}{\mathbb E}
\newcommand{\one}{\mathbf 1}
\newcommand{\Exp}{\operatorname{Exp}}
\newcommand{\dd}{\,d}

\title{Poisson--Dirichlet Universality for\\Luce Permutations}
\author{Christopher D. Long}
\date{Mechanical extraction --- audit pending}

\begin{document}
\maketitle

\begin{abstract}
Let a Luce permutation be generated by ranking independent exponential clocks
with prescribed positive rates.  Under a uniform bound on the ratio of the
largest to the smallest rate, the normalized ranked cycle lengths converge in
$\ell^1$ to $\operatorname{PD}(1)$, and their size-biased deletion order
converges in $\ell^1$ to $\operatorname{GEM}(1)$.  In particular, the cycle
length of a uniformly chosen label, divided by the number of labels, converges
to the uniform distribution on $(0,1)$.

This manuscript is a mechanical extraction of the macroscopic-cycle section of
the archived 139-page Luce manuscript at pinned repository commit
\texttt{6178a99}.  The theorem hypotheses,
conclusions, topology, labels, and proof content are preserved.  The extraction
does not constitute a new proof audit or an upgrade of mathematical status.
\end{abstract}

\section{Model, scope, and proof dependencies}

Let $X_{n,i}\sim\Exp(\theta_{n,i})$, $1\le i\le n$, be independent clocks
with positive rates.  The label-to-position Luce permutation is
\[
  \sigma_n(i)=1+\#\{j:X_{n,j}<X_{n,i}\}.
\]
Multiplying all rates by a common positive constant does not change the
permutation.

All model-specific ingredients used in the macroscopic-cycle proof are stated
and proved below: the subsequential marked kernel and its contraction, the
clock-slab and residual permanental structure, hereditary flatness, the
relative-prefix estimate, two-scale orbit balance, the rooted closure law,
delete-and-regenerate stability, and the finite-dimensional GEM induction.
The proof does not import the mesoscopic-cycle theorem, the total-cycle CLT, or
any result from another active manuscript in this repository.

The external inputs are standard compactness and disintegration facts for
probability measures on compact standard Borel spaces, Kolmogorov's maximal
inequality, Doob's $L^2$ maximal inequality, the continuous mapping theorem,
and the standard relation between $\operatorname{GEM}(1)$ and
$\operatorname{PD}(1)$; for the last point see Pitman
\cite[Chapter~3]{Pitman2006}.  The old-label/new-label register accompanying
this extraction records every theorem, proposition, lemma, equation label, and
bibliographic key moved from the archive.

\section{Macroscopic cycles under uniformly comparable weights}
\label{sec:macroscopic-cycles}

The preceding results concern cycles whose lengths remain fixed as $n\to\infty$.
We now consider cycles of order $n$.  Under a uniform bound on the ratio of the
largest to the smallest weight, the macroscopic cycle partition is universal.
We write $\operatorname{GEM}(1)$ for the residual-allocation sequence
\[
 U_1,\quad (1-U_1)U_2,\quad
 (1-U_1)(1-U_2)U_3,\ldots,
\]
where the $U_j$ are independent uniform random variables, and
$\operatorname{PD}(1)$ for its decreasing rearrangement; see
Pitman \cite[Chapter~3]{Pitman2006}.  Let
\[
 \mathcal S^\downarrow
 =\left\{x_1\ge x_2\ge\cdots\ge0:
 \sum_{j\ge1}x_j=1\right\},
\]
equipped with the $\ell^1$ metric.

\begin{theorem}[Macroscopic cycle universality]
\label{thm:macroscopic-pd}
Let $\sigma_n$ be a Luce permutation with positive weights
$\theta_{n,1},\ldots,\theta_{n,n}$ and assume
\[
 \sup_n\frac{\max_i\theta_{n,i}}{\min_i\theta_{n,i}}<\infty.
\]
If $L_{n,1}\ge L_{n,2}\ge\cdots$ are the cycle lengths of $\sigma_n$, then
\[
 \left(\frac{L_{n,1}}n,\frac{L_{n,2}}n,\ldots\right)
 \Rightarrow\operatorname{PD}(1)
\]
in $\mathcal S^\downarrow$.  Moreover, if
$\widetilde L_{n,1},\widetilde L_{n,2},\ldots$ are obtained by choosing a
uniform vertex, listing the length of its cycle, deleting that cycle, and
repeating on the residual vertex set, then
\[
 \left(\frac{\widetilde L_{n,1}}n,
 \frac{\widetilde L_{n,2}}n,\ldots\right)
 \Rightarrow\operatorname{GEM}(1)
\]
in $\ell^1$.  In particular, if $I_n$ is a uniformly chosen label and
$\mathscr L_n(I_n)$ is the length of its cycle, then
\[
        \frac{\mathscr L_n(I_n)}n\Rightarrow\operatorname{Unif}(0,1).
\]
\end{theorem}

The proof uses one fixed finite partition of clock time.  Conditional on the
clock slabs, each slab row is a permanental matching whose entries are nearly
constant.  A common minorization gives a spectral gap depending only on the
weight-comparability constant.  A relative-prefix theorem then controls every
target category under arbitrary adaptive exploration.  We formulate all
concentration statements uniformly over an explicit class of regular residual
histories, and prove that deleting an interior macroscopic cycle preserves that
class.  This supplies the regeneration needed for the $\operatorname{GEM}(1)$
induction.

Multiplying all weights by a common positive factor does not change the Luce
order.  We therefore normalize throughout the proof so that, for a fixed
$\kappa<\infty$,
\begin{equation}
        1\le\theta_{n,i}\le\kappa,
        \qquad 1\le i\le n.
        \tag{6.0}\label{eq:pd-normalized-rates}
\end{equation}

\subsection{Subsequential marked kernels and a uniform contraction}

Fix an arbitrary subsequence.  Passing to a further subsequence, compactness
gives
\[
 \eta_n:=\frac1n\sum_{i=1}^n\delta_{(i/n,\theta_{n,i})}
 \Longrightarrow\eta
\]
on $[0,1]\times[1,\kappa]$.  The first marginal of $\eta$ is Lebesgue
measure.  Choose a disintegration
\[
        \eta(\dd x,\dd\vartheta)=\dd x\,\eta_x(\dd\vartheta),
\]
defined for almost every $x$.  Put
\[
 A(t)=\int(1-e^{-\vartheta t})\,\eta(\dd x,\dd\vartheta),
 \qquad
 B(t)=A'(t)=\int\vartheta e^{-\vartheta t}\,
 \eta(\dd x,\dd\vartheta),
\]
let $\tau=A^{-1}$, and define, for almost every $x$,
\[
 g_x(t)=\int\vartheta e^{-\vartheta t}\,\eta_x(\dd\vartheta),
 \qquad
 \rho(x,y)=\frac{g_x(\tau(y))}{B(\tau(y))}.
\]
Let $K$ be the associated Markov operator on $L^2[0,1]$.

\begin{lemma}[Bistochasticity and common minorization]
\label{lem:pd-kernel-contraction}
The kernel $\rho$ is bistochastic up to null sets and
\[
        \iint_{[0,1]^2}\rho(x,y)^2\dd x\dd y\le\kappa.
\]
On the mean-zero subspace,
\[
        \|K\|_{L^2_0\to L^2_0}
        \le q_\kappa:=\sqrt{1-\kappa^{-1}}<1,
\]
and the same bound holds for $K^*$.
\end{lemma}

\begin{proof}
Using $y=A(t)$ and $\dd y=B(t)\dd t$,
\[
 \int_0^1\rho(x,y)\dd y=\int_0^\infty g_x(t)\dd t=1
\]
for almost every $x$, while
\[
 \int_0^1\rho(x,y)\dd x
 =\frac{\int_0^1g_x(\tau(y))\dd x}{B(\tau(y))}=1.
\]
Jensen's inequality and $\vartheta\le\kappa$ give
\[
 \int_0^1g_x(t)^2\dd x
 \le\int\vartheta^2e^{-2\vartheta t}\,
 \eta(\dd x,\dd\vartheta)
 \le\kappa B(2t).
\]
Since $\vartheta\ge1$, $B(2t)\le e^{-t}B(t)$, and hence
\[
 \iint\rho^2
 =\int_0^\infty\frac{\int g_x(t)^2\dd x}{B(t)}\dd t
 \le\kappa.
\]

For the contraction, regard $g_x(t)\dd t$ as the conditional law of a clock
time $T$ given $x$.  From \eqref{eq:pd-normalized-rates},
$g_x(t)\ge e^{-\kappa t}$.  Let
\[
        \alpha=\kappa^{-1},
        \qquad \nu(\dd t)=\kappa e^{-\kappa t}\dd t.
\]
There are probability kernels $R_x$ such that
\[
        g_x(t)\dd t=\alpha\nu(\dd t)+(1-\alpha)R_x(\dd t).
\]
The marginal law of $T$ is
$\mu=\alpha\nu+(1-\alpha)\bar R$, where
$\bar R=\int_0^1R_x\dd x$.  If $f\in L^2(\mu)$ and $\mu(f)=0$, then
\[
 \int f(t)g_x(t)\dd t
 =(1-\alpha)\{R_x(f)-\bar R(f)\}.
\]
Jensen's inequality therefore yields
\[
\begin{aligned}
 \int_0^1\left|\int f(t)g_x(t)\dd t\right|^2\dd x
 &\le(1-\alpha)^2\bar R(f^2)\\
 &\le(1-\alpha)\mu(f^2).
\end{aligned}
\]
The change of variables $y=A(t)$ is an isometry from $L^2(\dd y)$ to
$L^2(\mu)$, and under this isometry the displayed conditional-expectation
operator is $K$.  This proves the bound; the adjoint has the same norm.
\end{proof}

\subsection{Clock slabs and residual permanental matchings}

Fix $T>0$ and a partition
\[
        0=t_0<t_1<\cdots<t_d=T,
        \qquad \max_a(t_a-t_{a-1})\le h,
\]
and append the tail interval $I_\star=(T,\infty)$.  Put
\[
 I_a=(t_{a-1},t_a]\quad(a\le d),
 \qquad H_{n,a}=\{i:X_{n,i}\in I_a\},
 \qquad m_{n,a}=|H_{n,a}|.
\]
If $M_{n,a}=\sum_{c\le a}m_{n,c}$, then the ranks occupied by $H_{n,a}$ are
\[
        Q_{n,a}=\{M_{n,a-1}+1,\ldots,M_{n,a}\};
\]
the tail row is defined analogously.  Set
\[
 N_{n,ab}=|Q_{n,a}\cap H_{n,b}|,
 \qquad p_{n,a}=m_{n,a}/n.
\]
When $m_{n,a}>0$, put $P_{n,ab}=N_{n,ab}/m_{n,a}$; if
$m_{n,a}=0$, define that row of $P_n$ arbitrarily.  On the good events used
below every row is nonempty.  Thus $P_n$ is stochastic and $p_nP_n=p_n$.

Let
\[
 J_a=(A(t_{a-1}),A(t_a)]\quad(a\le d),
 \qquad J_\star=(A(T),1],
 \qquad p_a=|J_a|,
\]
and define
\[
 q_{ab}=\int_{J_a\times J_b}\rho(x,y)\dd x\dd y,
 \qquad P_{ab}=q_{ab}/p_a.
\]
Every $p_a$ and $q_{ab}$ is strictly positive.

\begin{lemma}[Slab law and deterministic good events]
\label{lem:pd-slab-law}
For every fixed partition,
\[
        p_{n,a}\xrightarrow{\Pp}p_a,
        \qquad N_{n,ab}/n\xrightarrow{\Pp}q_{ab}
\]
simultaneously in $a,b$.  There are deterministic $\varepsilon_n\downarrow0$
and events $\mathcal G_n$ with $\Pp(\mathcal G_n)\to1$ such that, on
$\mathcal G_n$:
\begin{enumerate}[label=(\roman*)]
\item all of the preceding convergences hold with error at most
$\varepsilon_n^2$, and, for all sufficiently large $n$,
\[
        p_{n,a}\ge\tfrac12p_a,
        \qquad N_{n,ab}\ge\tfrac12q_{ab}n
\]
for every $a,b$;
\item for a constant $q_0\in(q_\kappa,1)$,
\[
        \|P_n\|_{L^2_0(p_n)\to L^2_0(p_n)}\le q_0;
\]
\item if $S$ is a partition endpoint and $J_b\subset[0,A(S)]$, then
\[
        P_{n,ab}\le 2M_Sp_{n,b}\qquad\text{for every }a,
        \qquad M_S=\kappa e^{(\kappa-1)S};
        \tag{6.1}
\]
\item for every union $D$ of slab states,
\[
 \sum_{a\in D}\sum_b p_{n,a}\frac{P_{n,ab}^2}{p_{n,b}}
 \le2\kappa;
        \tag{6.2}
\]
\item if $U$ is a partition endpoint and $D(U)$ is the union of rows whose
clock intervals lie in $(U,\infty)$, then
\[
        \sum_{a\in D(U)}p_{n,a}\le2e^{-U}+\varepsilon_n^2.
        \tag{6.3}
\]
\end{enumerate}
The constants may depend on the fixed partition and on $S$, but not on $n$.
\end{lemma}

\begin{proof}
For a fixed clock cell $I_b=(u,v]$, the variables
$\one_{\{X_{n,i}\in I_b\}}$ are independent and have means
$e^{-\theta_{n,i}u}-e^{-\theta_{n,i}v}$.  Kolmogorov's maximal inequality
gives
\[
 \max_{k\le n}\left|\sum_{i\le k}
 \left(\one_{\{X_{n,i}\in I_b\}}
 -\Ee\one_{\{X_{n,i}\in I_b\}}\right)\right|=o_{\Pp}(n).
\]
Marked empirical convergence and the Lebesgue first marginal imply, uniformly
in $x\in[0,1]$,
\[
 \frac1n\sum_{i\le nx}
 \Ee\one_{\{X_{n,i}\in I_b\}}
 \longrightarrow
 \int_0^x\int
 (e^{-\vartheta u}-e^{-\vartheta v})
 \eta_z(\dd\vartheta)\dd z.
\]
This proves the slab-size convergence and convergence of the endpoints of
$Q_{n,a}/n$ to those of $J_a$.  Applying the same uniform partial-sum estimate
at those random endpoints gives
\[
 \frac{N_{n,ab}}n\longrightarrow
 \int_{J_a}\int
 (e^{-\vartheta t_{b-1}}-e^{-\vartheta t_b})
 \eta_x(\dd\vartheta)\dd x=q_{ab}.
\]
The last equality follows by integrating $\rho(x,y)$ over $J_b$ and using
$y=A(t)$.

The matrix $P$ is the compression of $K$ to functions that are constant on the
sets $J_a$: explicitly, it is $\Pi_{\mathcal J}K\Pi_{\mathcal J}$.  Conditional
expectation is a contraction, so Lemma \ref{lem:pd-kernel-contraction} gives
$\|P\|_{L^2_0(p)}\le q_\kappa$.  Entrywise convergence of the conjugated
matrices
$D_{p_n}^{1/2}P_nD_{p_n}^{-1/2}$ proves (ii).

If $y\le A(S)$, then $\tau(y)\le S$ and
\[
 \rho(x,y)
 =\frac{g_x(\tau(y))}{B(\tau(y))}
 \le\frac{\kappa e^{-\tau(y)}}{e^{-\kappa\tau(y)}}
 \le\kappa e^{(\kappa-1)S}=:M_S.
\]
Hence $P_{ab}\le M_Sp_b$, which yields (6.1) by entrywise convergence.  On
every rectangle, Cauchy's inequality gives
\[
 p_a\frac{P_{ab}^2}{p_b}
 =\frac{(\int_{J_a\times J_b}\rho)^2}{p_ap_b}
 \le\int_{J_a\times J_b}\rho^2.
\]
Summing and using Lemma \ref{lem:pd-kernel-contraction} proves the limiting
version of (6.2), and entrywise convergence gives its finite-$n$ form.  Finally,
the expected proportion of clocks exceeding $U$ is at most $e^{-U}$ because
all rates are at least one; the slab law gives (6.3).  Since the partition is finite, choose deterministic $\delta_n\downarrow0$
so that all convergence assertions hold with error at most $\delta_n$ on a
single event of probability tending to one, and set
$\varepsilon_n=\sqrt{\delta_n}$.  This gives the stated event
$\mathcal G_n$ and the errors $\varepsilon_n^2$ in (i) and (6.3).
\end{proof}

Condition on the complete slab data: the sets $H_{n,a}$, the rank intervals
$Q_{n,a}$, and the ordered clock values
\[
        s_{a,1}<\cdots<s_{a,m_{n,a}}
\]
in every slab.  Write
$Q_{n,a}=\{q_{a,1}<\cdots<q_{a,m_{n,a}}\}$, and let
$\phi_a:H_{n,a}\to[m_{n,a}]$ be defined by
$\sigma_n(i)=q_{a,\phi_a(i)}$.  Then
\begin{equation}
 \Pp(\phi_a=\phi\mid\text{slab data})
 =\frac{\prod_{i\in H_{n,a}}e^{-\theta_{n,i}s_{a,\phi(i)}}}
 {\operatorname{per}(e^{-\theta_{n,i}s_{a,j}})_{i,j}}.
 \tag{6.3a}\label{eq:pd-permanental-row}
\end{equation}
The row matchings are conditionally independent.  The factors
$\prod_i\theta_{n,i}$ cancel from the conditional density.

\begin{lemma}[Residual permanental factorization]
\label{lem:pd-residual-factorization}
Condition on the slab data and on any compatible partial collection of exposed
matching edges.  In every row, delete the exposed source and target vertices.
The remaining row matchings are conditionally independent, and each has the
permanental law \eqref{eq:pd-permanental-row} on its remaining submatrix.
\end{lemma}

\begin{proof}
The conditional density of the complete matching is a product over rows and,
within each row, a product over its edges.  Conditioning on fixed exposed edges
removes their constant factors and restricts the sum to bijections of the
remaining rows and columns.  The normalizing sum is the permanent of the
remaining submatrix, and the factorization over rows is unchanged.
\end{proof}

\begin{lemma}[Hereditary flatness]
\label{lem:pd-hereditary-flatness}
Let $A=(a_{ij})$ be a positive $m\times m$ matrix satisfying
\[
        e^{-\epsilon}\le a_{ij}/a_{ik}\le e^\epsilon
\]
for every row $i$ and every pair of columns $j,k$.  Under the permanental
matching law proportional to $\prod_i a_{i,\phi(i)}$, after conditioning on an
arbitrary partial matching, every remaining edge probability satisfies
\[
 \frac{e^{-2\epsilon}}s
 \le\Pp(\phi(i)=j\mid\text{partial matching})
 \le\frac{e^{2\epsilon}}s,
\]
where $s$ rows and columns remain.  For a Luce slab of width at most $h$, this
holds with $\epsilon=\kappa h$; write
\[
        \beta=2\kappa h.
        \tag{6.4}
\]
\end{lemma}

\begin{proof}
It suffices to work with the current remaining submatrix.  For two remaining
columns $j,k$,
\[
 \frac{\Pp(\phi(i)=j)}{\Pp(\phi(i)=k)}
 =\frac{a_{ij}\operatorname{per}A_{-i,-j}}
 {a_{ik}\operatorname{per}A_{-i,-k}}.
\]
Map a term of the first minor to a term of the second by locating the row sent
to $k$ and sending that row to $j$ instead.  Corresponding products differ by
a factor in $[e^{-\epsilon},e^\epsilon]$.  Including the factor
$a_{ij}/a_{ik}$ shows that every pair of edge probabilities has ratio in
$[e^{-2\epsilon},e^{2\epsilon}]$.  Since the $s$ probabilities sum to one,
the individual bounds follow.  In a clock slab,
\[
        e^{-\kappa h}\le
        e^{-\theta_i(s_j-s_k)}\le e^{\kappa h}.
\]
\end{proof}
\begin{lemma}[Predictable row exposure]
\label{lem:pd-predictable-exposure}
Condition on the slab data and on a compatible partial collection of exposed
matching edges.  Let $(\mathscr F_r)$ be an exploration filtration such that,
before the $(r+1)$st edge is exposed, its source is an unexposed source that is
$\mathscr F_r$-measurable.  Then, conditional on $\mathscr F_r$, the target of
that source has the one-edge marginal of the residual permanental matching in
its row.  In particular, if that row has $s$ targets remaining and has width at
most $h$, every individual target has conditional probability in
\[
        [e^{-\beta}/s,e^\beta/s],
        \qquad \beta=2\kappa h.
\]
If a distinguished target $v$ in that row is forbidden and one conditions the
next target not to equal $v$, then every one of the other $s-1$ targets has
conditional probability in
\[
        [e^{-\beta}/(s-1),e^\beta/(s-1)].
\]
The conclusions remain true when the next source is selected adaptively using
all previously exposed edges, including edges from the same row.
\end{lemma}

\begin{proof}
By Lemma \ref{lem:pd-residual-factorization}, conditional on
$\mathscr F_r$ the unexposed part of the relevant row has the permanental law
on its current submatrix.  Lemma \ref{lem:pd-hereditary-flatness} therefore
gives the first interval of probabilities.  If $v$ is forbidden, conditioning
away from $v$ merely renormalizes the probabilities of the other targets.
Their pairwise ratios are unchanged and lie in $[e^{-\beta},e^\beta]$; since
their conditional probabilities sum to one, each lies in the second displayed
interval.  The measurability assumption is precisely what is needed to apply
the residual-factorization statement before the new edge is revealed.  No
independence between the chosen source and the previously exposed edges is
required.
\end{proof}


\subsection{A quantitative relative-prefix theorem}

\begin{lemma}[Relative prefix law]
\label{lem:pd-relative-prefix}
Fix $d\ge1$ and $c_*>0$.  There are absolute constants
$\beta_0>0$ and $C<\infty$, and deterministic sequences
$\chi_m\downarrow0$ and $\pi_m\downarrow0$, depending only on $d$ and $c_*$,
with the following property.  Let $m$ targets be
partitioned into $d$ categories of sizes $m_1,\ldots,m_d$ satisfying
$m_b\ge c_*m$.  Targets are selected sequentially without replacement in an
arbitrary predictable, adaptive source order.  Suppose that, whenever $s$
targets remain, every individual remaining target has conditional selection
probability in
\[
        [e^{-\beta}/s,e^\beta/s],
        \qquad 0\le\beta\le\beta_0.
\]
If $X_b(r)$ is the number of category-$b$ targets among the first $r$
selections, then
\[
 \Pp\left(
 \max_{b\le d}\max_{0\le r\le m}
 \frac{|X_b(r)-rm_b/m|}{m_b}>C\beta+\chi_m
 \right)\le\pi_m.
 \tag{6.5}\label{eq:pd-relative-prefix}
\]
The estimate is uniform over all adaptive source rules satisfying the displayed
one-target bounds.  It remains valid after deleting any fixed number of
distinguished targets before the procedure, or after forbidding those targets
throughout the procedure, after changing $c_*$ and the two deterministic
sequences harmlessly.
\end{lemma}

\begin{proof}
Fix one category and write
\[
        N_r=m_b-X_b(r),\qquad s_r=m-r.
\]
Let $Y_{r+1}$ be the indicator that the next selected target belongs to this
category.  By summing the individual-target bounds over the $N_r$ remaining
targets in the category,
\[
 \pi_r:=\Ee(Y_{r+1}\mid\mathscr F_r)
 \in\left[e^{-\beta}\frac{N_r}{s_r},
          e^\beta\frac{N_r}{s_r}\right].
 \tag{6.6}\label{eq:pd-prefix-pi}
\]
Set $a_m=m_b^{2/3}$ and stop at
\[
        \tau=\inf\{r:N_r<a_m\text{ or }s_r<a_m\}.
\]
For $r<\tau$, put $Z_r=\log(N_r/s_r)$.  Since
\[
 Z_{r+1}-Z_r
 =-\log(1-s_r^{-1})
  +Y_{r+1}\log(1-N_r^{-1}),
\]
Taylor's formula and \eqref{eq:pd-prefix-pi} give
\[
 \left|\Ee(Z_{r+1}-Z_r\mid\mathscr F_r)\right|
 \le C_1\frac\beta{s_r}
 +C_1\left(\frac1{s_r^2}+\frac1{s_rN_r}\right),
 \tag{6.7}\label{eq:pd-prefix-drift}
\]
and the conditional variance of the centered increment is at most
$C_1/(s_rN_r)$.  Because $m_b\ge c_*m$,
\[
 \sum_{r<\tau}\frac1{s_rN_r}
 \le a_m^{-1}\sum_{s\ge a_m}\frac1s
 \le C_{c_*}\frac{1+\log m}{m^{2/3}}=:v_m,
 \tag{6.8}\label{eq:pd-prefix-qv}
\]
and the sum of the two second-order drift terms in
\eqref{eq:pd-prefix-drift} is $O(v_m)$.

Let $M_r$ be the stopped martingale obtained by summing the centered
increments.  Doob's $L^2$ inequality and \eqref{eq:pd-prefix-qv} imply
\[
 \Pp\left(\max_{r\le\tau}|M_r|>v_m^{1/4}\right)
 \le C v_m^{1/2}.
 \tag{6.9}\label{eq:pd-prefix-doob}
\]
On the complementary event, for every $r<\tau$,
\[
 \left|\log\frac{N_r/m_b}{s_r/m}\right|
 \le C_2\beta\log\frac m{s_r}+C_2(v_m^{1/4}+v_m).
 \tag{6.10}\label{eq:pd-prefix-log}
\]
Choose $\beta_0$ so that $C_2\beta_0\le1/2$.  Writing $u=s_r/m$ and
$\eta_m=C_2(v_m^{1/4}+v_m)$, \eqref{eq:pd-prefix-log} yields
\[
 e^{-\eta_m}u^{1+C_2\beta}
 \le\frac{N_r}{m_b}
 \le e^{\eta_m}u^{1-C_2\beta}.
 \tag{6.11}\label{eq:pd-prefix-power}
\]
For $|a|\le1/2$,
\[
        \sup_{0\le u\le1}|u^{1+a}-u|\le C|a|,
\]
which follows from the mean-value theorem in the exponent and
$\sup_{0<u\le1}u|\log u|<\infty$.  Hence, before $\tau$,
\[
        \left|\frac{N_r}{m_b}-\frac{s_r}{m}\right|
        \le C_3\beta+C_3\eta_m.
 \tag{6.12}\label{eq:pd-prefix-before-stop}
\]

At the stopping time, if $s_\tau<a_m$, then
$s_\tau/m=O(m^{-1/3})$ and the upper bound in
\eqref{eq:pd-prefix-power}, applied one step earlier, gives
$N_\tau/m_b=o(1)$.  If $N_\tau<a_m$, then the lower bound in
\eqref{eq:pd-prefix-power} gives $s_\tau/m=o(1)$.  In either case, after
$\tau$ both normalized remaining counts are nonincreasing and their difference
is bounded by $C_3\beta+o(1)$.  Thus
\eqref{eq:pd-prefix-before-stop} extends to all $0\le r\le m$.

For example, one may take
\[
 \chi_m=C_{d,c_*}\left(v_m^{1/4}+m^{-1/12}\right),
 \qquad
 \pi_m=C_{d,c_*}v_m^{1/2}.
\]
A union bound over the fixed number of categories proves
\eqref{eq:pd-relative-prefix}.  Deleting or forbidding finitely many targets
changes every relevant category size and the total size by $O(1)$; for all
large $m$ the positive-proportion hypothesis and the same one-target ratio
bounds remain valid after changing $c_*$ by a fixed factor.
\end{proof}

In every application below the mesh is chosen sufficiently fine that
$\beta=2\kappa h\le\beta_0$.

\subsection{Regular residual histories and scale-sensitive orbit balance}

A residual history $\mathcal H$ consists of the complete slab data, all roots
selected so far, and the matching edges in a collection of fully exposed and
deleted cycles; it contains no additional information about unexposed row
matchings.  After such a history, let $R\subset[n]$ be the remaining vertex set
and $N=|R|$.  Put
\[
 m_a^R=|H_{n,a}\cap R|,
 \qquad N_{ab}^R=|Q_{n,a}\cap H_{n,b}\cap R|,
 \qquad r=N/n.
\]
For every deleted cycle $C$ and every slab $a$,
\begin{equation}
        |C\cap H_{n,a}|=|C\cap Q_{n,a}|,
        \label{eq:pd-cycle-block-equality}
\end{equation}
because $\sigma_n$ maps $C\cap H_{n,a}$ bijectively onto
$C\cap Q_{n,a}$.  Consequently $m_a^R=|Q_{n,a}\cap R|$.

Let $U$ be a master-partition endpoint.  A residual history is called
$(U,\delta,\xi)$-regular if $N\ge\xi n$, $0\le\delta<1/4$, and, for every row
whose clock interval is contained in $[0,U]$ and every target category $b$,
\begin{equation}
 \left|\frac{m_a^R}{r m_{n,a}}-1\right|\le\delta,
 \qquad
 \left|\frac{N_{ab}^R}{r N_{n,ab}}-1\right|\le\delta.
 \label{eq:pd-regular-history}
\end{equation}
The initial history is $(U,0,1)$-regular.  By Lemma
\ref{lem:pd-residual-factorization}, conditional on any residual history, the
unexposed row matchings are independent permanental matchings on their
remaining submatrices.

An \emph{admissible exploration} starts from a residual history and reveals one
outgoing matching edge at a time.  Before an edge is revealed, its source must
be measurable with respect to the filtration generated by the slab data, the
residual history, and all previously exposed edges.  We also permit a fixed
number of distinguished targets to be forbidden, with the next edge in a row
containing such a target sampled conditionally on avoiding it.  Lemma
\ref{lem:pd-predictable-exposure} shows that every actual cycle exploration and
every root-avoiding exploration used below is admissible.

Consider an injective orbit segment $v_0,v_1,\ldots,v_\ell$, where
$v_{j+1}=\sigma_n(v_j)$ for the actual exploration, and allow
$v_\ell=v_0$ only at a terminal closing step.  For an auxiliary exploration,
$v_{j+1}$ denotes the target selected under its conditional avoid-target law.
Put
\[
 R_a(\ell)=\#\{0\le j<\ell:v_j\in H_{n,a}\},
\]
and
\[
 R_{ab}(\ell)=\#\{0\le j<\ell:v_j\in H_{n,a},\ v_{j+1}\in H_{n,b}\}.
\]
Whenever no ambiguity is possible we suppress the argument $\ell$.  The
deterministic flow identity is
\begin{equation}
 R_b-\sum_aR_{ab}
 =\one_{\{v_0\in H_{n,b}\}}-
  \one_{\{v_\ell\in H_{n,b}\}}.
 \label{eq:pd-flow}
\end{equation}

\begin{lemma}[Scale-sensitive two-scale orbit balance]
\label{lem:pd-two-scale-balance}
Fix master-partition endpoints $0<S<U\le T$, $\xi>0$, and
$0\le\delta<1/4$.  There are deterministic sequences
$\chi_n\downarrow0$ and $\pi_n\downarrow0$, depending on the fixed partition,
$U$, and $\xi$ but not on the residual history, with the following property.
On $\mathcal G_n$, uniformly over every $(U,\delta,\xi)$-regular history and
every admissible exploration, with conditional probability at least
$1-\pi_n$ the following assertions hold simultaneously at every defined prefix
length $\ell\ge1$.

For every row $a$ whose clock interval is contained in $[0,U]$ and every target
category $b$,
\begin{equation}
 \left|R_{ab}(\ell)-\frac{R_a(\ell)}{m_a^R}N_{ab}^R\right|
 \le (C\beta+\chi_n)N_{ab}^R,
 \label{eq:pd-row-prefix}
\end{equation}
where $C$ is absolute and $\beta=2\kappa h$.  Define
\[
        \Delta_n(U,\delta,h)
        =\delta+\beta+\chi_n+\varepsilon_n+e^{-U/2}.
\]
Then, for every inner state $b$ with $J_b\subset[0,A(S)]$,
\begin{equation}
 \left|\frac{R_b(\ell)}{\ell p_{n,b}}-1\right|
 \le C_{\kappa,S}\frac n\ell\,
        \Delta_n(U,\delta,h)
      +\frac{C_{\mathcal J}}\ell.
 \label{eq:pd-scale-sensitive-balance}
\end{equation}
Here $C_{\kappa,S}$ is independent of the mesh and of $\ell$; the finite
constant $C_{\mathcal J}$ may depend on the fixed master partition.  In
particular, for every fixed $\lambda>0$, all prefixes with $\ell\ge\lambda n$
satisfy
\[
 \max_{b:J_b\subset[0,A(S)]}
 \left|\frac{R_b(\ell)}{\ell p_{n,b}}-1\right|
 \le C_{\kappa,S,\lambda}\Delta_n(U,\delta,h)+o(1).
\]
\end{lemma}

\begin{proof}
On $\mathcal G_n$, every $m_{n,a}$ and every $N_{n,ab}$ is bounded below by a
positive constant times $n$.  If the history is regular, the same is true for
$m_a^R$ and $N_{ab}^R$, uniformly over the stated histories.  Consequently,
for each row below $U$, the residual target categories have proportions bounded
below by a positive constant depending only on the fixed partition.  Apply
Lemmas \ref{lem:pd-predictable-exposure} and
\ref{lem:pd-relative-prefix} to the adaptive sequence of queried sources in
that row, extending it predictably after the orbit segment ends until all
ordinary targets in that row have been selected.  In a row with a forbidden
target, this means a complete order of the allowable targets.  A union bound over the finitely many rows and categories gives
deterministic sequences $\chi_n,\pi_n\downarrow0$ for which
\eqref{eq:pd-row-prefix} holds simultaneously for every prefix.  If finitely
many distinguished targets are forbidden, remove them from the corresponding
category counts; the resulting $O(1)$ changes are absorbed into $\chi_n$.

Write $p=p_n$, $P=P_n$, $q_a=R_a/\ell$, and
\[
        E_{ab}=R_{ab}-R_aP_{ab}.
\]
For a row below $U$, regularity gives
\[
 \frac{N_{ab}^R}{m_a^R}=P_{ab}\{1+O(\delta)\},
\]
uniformly, and $R_a\le m_a^R$.  Hence
\begin{equation}
        |E_{ab}|
        \le C(\delta+\beta+\chi_n)N_{n,ab}
        \qquad(a\notin D(U)),
 \label{eq:pd-good-row-error}
\end{equation}
where $D(U)$ is the union of rows whose clock intervals lie in $(U,\infty)$.
For an outer row we use only injectivity of the explored path:
$R_{ab}\le N_{ab}^R\le N_{n,ab}$ and
$R_aP_{ab}\le m_a^RP_{ab}\le N_{n,ab}$.  Thus
\begin{equation}
        |E_{ab}|\le2N_{n,ab}
        \qquad(a\in D(U)).
 \label{eq:pd-bad-row-error}
\end{equation}

Let
\[
        G_b=\sum_{a\notin D(U)}E_{ab},
        \qquad H_b=\sum_{a\in D(U)}E_{ab}.
\]
Since $\sum_aN_{n,ab}=m_{n,b}=np_b$, equation
\eqref{eq:pd-good-row-error} gives
\begin{equation}
 \left\|\frac{G}{\ell p}\right\|_{L^2(p)}
 \le C\frac n\ell(\delta+\beta+\chi_n).
 \label{eq:pd-good-L2}
\end{equation}
For the outer rows, Cauchy's inequality and the Hilbert--Schmidt block estimate
on $\mathcal G_n$ give
\begin{align}
 \sum_b\frac{H_b^2}{m_{n,b}}
 &\le4\sum_b\frac{(\sum_{a\in D(U)}N_{n,ab})^2}{m_{n,b}}\notag\\
 &\le4m_{n,D(U)}
 \sum_{a\in D(U)}\sum_b
 \frac{m_{n,a}P_{n,ab}^2}{m_{n,b}}
 \le C m_{n,D(U)}.
 \label{eq:pd-bad-L2-raw}
\end{align}
Therefore, using the tail-mass estimate on $\mathcal G_n$,
\begin{equation}
 \left\|\frac{H}{\ell p}\right\|_{L^2(p)}
 \le C\frac n\ell
       \left(e^{-U/2}+\varepsilon_n\right).
 \label{eq:pd-bad-L2}
\end{equation}

Let $d_b$ denote the right-hand side of \eqref{eq:pd-flow}.  It has at most
two nonzero coordinates, each of absolute value at most one.  Dividing
\eqref{eq:pd-flow} by $\ell$ gives
\[
        q-qP=e,
        \qquad e_b=\frac{d_b+G_b+H_b}{\ell}.
\]
Since the partition is fixed,
\begin{equation}
 \left\|\frac e p\right\|_{L^2(p)}
 \le C\frac n\ell\Delta_n(U,\delta,h)
      +\frac{C_{\mathcal J}}\ell.
 \label{eq:pd-g-L2-scale}
\end{equation}
Put $u_a=q_a/p_a-1$.  Because $pP=p$ and $\sum_ap_au_a=0$,
\[
        u-P^*u=e/p.
\]
The spectral-gap event in $\mathcal G_n$ yields
\begin{equation}
 \|u\|_{L^2(p)}
 \le(1-q_0)^{-1}\left\|\frac e p\right\|_{L^2(p)}.
 \label{eq:pd-u-L2-scale}
\end{equation}

It remains to recover a coordinatewise estimate for an inner state $b$.
The inner-column bound and stationarity imply
\[
 \sum_a p_a\left(\frac{P_{ab}}{p_b}\right)^2
 \le \left(\max_a\frac{P_{ab}}{p_b}\right)
      \sum_a\frac{p_aP_{ab}}{p_b}
 \le C_{\kappa,S}.
\]
Hence
\[
        |(P^*u)_b|\le C_{\kappa,S}\|u\|_{L^2(p)}.
\]
For the coordinate $(e/p)_b$, equations
\eqref{eq:pd-good-row-error} and \eqref{eq:pd-bad-row-error}, together with
$P_{ab}\le C_{\kappa,S}p_b$ for an inner column, give
\[
 \left|\frac{e_b}{p_b}\right|
 \le C_{\kappa,S}\frac n\ell
      \left(\delta+\beta+\chi_n+e^{-U}+\varepsilon_n^2\right)
      +\frac{C_{\mathcal J}}\ell.
\]
Combining this estimate with
$u_b=(e/p)_b+(P^*u)_b$, \eqref{eq:pd-g-L2-scale}, and
\eqref{eq:pd-u-L2-scale} proves
\eqref{eq:pd-scale-sensitive-balance}.  Every probabilistic input was the
conditional relative-prefix event, so the failure probability $\pi_n$ is
uniform over the entire stated class of histories and explorations.
\end{proof}

\subsection{The cycle of a uniform residual root}

For a vertex $v$, let $c(v)$ be the unique row index for which
$v\in Q_{n,c(v)}$.  An orbit can hit $v$ only when its current source lies in
$H_{n,c(v)}$.

\begin{lemma}[Root-avoiding change of measure]
\label{lem:pd-root-avoiding}
Fix a residual history and a distinguished remaining target $v$.  Explore the
path from $v$.  At the $j$th opportunity to hit $v$, let $\lambda_j$ be the
conditional probability, under the original residual matching law and given
that $v$ has not yet been selected, that the next target is $v$.  Let
$\Pp^\circ$ be the auxiliary law obtained by conditioning each such choice not
to select $v$.  If $S(k)$ is the number of closure opportunities among the
first $k$ completed edges, then
\begin{equation}
        \Pp(|\mathcal C(v)|>k\mid\mathcal H,v)
        =\Ee^\circ\prod_{j=0}^{S(k)-1}(1-\lambda_j),
 \label{eq:pd-root-rn}
\end{equation}
where the product is interpreted as zero if the auxiliary exploration is
forced to stop before completing $k$ edges.  If the closure row initially has
$m$ available targets, including $v$, then at its $j$th closure opportunity,
\begin{equation}
        \frac{e^{-\beta}}{m-j}\le\lambda_j\le
        \frac{e^\beta}{m-j},
        \qquad \beta=2\kappa h.
 \label{eq:pd-root-hazard}
\end{equation}
\end{lemma}

\begin{proof}
At a closure opportunity, the original transition law is the mixture of an
atom of mass $\lambda_j$ at $v$ and the conditional avoid-$v$ law of mass
$1-\lambda_j$.  A non-root previously visited vertex cannot be selected,
because its target column has already been used; thus the only possible first
repeat is the root.  Multiplying the sequential Radon--Nikodym factors gives
\eqref{eq:pd-root-rn}.  Before the $j$th opportunity, exactly $j$ non-$v$
targets have been removed from the closure row, so $m-j$ targets remain.
Equation \eqref{eq:pd-root-hazard} follows from hereditary flatness.  If only
$v$ remains, the next closure probability is one, and the survival product is
zero, which agrees with the stated convention.
\end{proof}

\begin{lemma}[Harmonic-product approximation]
\label{lem:pd-harmonic-product}
Fix $\zeta\in(0,1)$.  Uniformly for $c\in[1/2,2]$, integers $m\ge1$, and
$0\le s\le(1-\zeta)m$ for which all factors are nonnegative,
\[
 \prod_{j=0}^{s-1}\left(1-\frac c{m-j}\right)
 =(1-s/m)^c\left\{1+O_\zeta(m^{-1})\right\}.
\]
Moreover, uniformly for $x\in[\zeta,1-\zeta]$ and $|a|\le1/2$,
\[
        |(1-x)^{1+a}-(1-x)|\le C_\zeta|a|.
\]
\end{lemma}

\begin{proof}
Since $m-j\ge\zeta m$,
\[
 \sum_{j<s}\log\left(1-\frac c{m-j}\right)
 =-c\sum_{j<s}\frac1{m-j}+O_\zeta(m^{-1}).
\]
The harmonic sum equals
$\log(m/(m-s))+O_\zeta(m^{-1})$, which proves the first assertion after
exponentiation.  The second follows from the mean-value theorem in the exponent
and the boundedness of $|(1-x)^b\log(1-x)|$ for
$x\in[\zeta,1-\zeta]$ and $b\in[1/2,3/2]$.
\end{proof}

\begin{proposition}[Uniform conditional residual fraction]
\label{prop:pd-uniform-residual-fraction}
Fix master-partition endpoints $0<S<U\le T$, $\xi>0$,
$\zeta\in(0,1/2)$, and $0\le\delta<1/4$.  On $\mathcal G_n$, let
$\mathcal H$ be a $(U,\delta,\xi)$-regular history.  Choose a root uniformly
from its $N$ remaining vertices, independently of the unexposed matchings, and
let $L$ be the root-cycle length.  With $\chi_n,\pi_n$ as in Lemma
\ref{lem:pd-two-scale-balance}, uniformly over all such histories,
\begin{align}
 &\sup_{x\in[\zeta,1-\zeta]}
 \left|\Pp\left(\left.\frac LN\le x\right|\mathcal H\right)-x\right|
 \notag\\
 &\qquad\le
 C_{\kappa,S,\xi,\zeta}
 \bigl(\delta+\beta+\chi_n+\varepsilon_n+e^{-U/2}\bigr)
 +C_\xi(e^{-S}+\varepsilon_n^2)
 +\pi_n+\frac{C_{\mathcal J,\xi,\zeta}}n.
 \label{eq:pd-root-cdf}
\end{align}
The constants multiplying the principal error are independent of the mesh;
$C_{\mathcal J,\xi,\zeta}$ may depend on the fixed master partition.
\end{proposition}

\begin{proof}
Let $v$ be the uniform root and put $C=c(v)$.  Complete-cycle deletion gives
$|Q_{n,a}\cap R|=m_a^R$, so, conditional on the history,
\[
        \Pp(C=a\mid\mathcal H)=m_a^R/N.
\]
Consequently, using $N\ge\xi n$ and the tail bound on $\mathcal G_n$,
\begin{equation}
 \Pp(C\in D(S)\mid\mathcal H)
 \le C_\xi(e^{-S}+\varepsilon_n^2).
 \label{eq:pd-closure-tail}
\end{equation}
We henceforth work on the event $C\notin D(S)$.  Put
$m=m_C^R$ and $p_C=p_{n,C}$.  Regularity gives
\begin{equation}
        \frac{m}{Np_C}=1+O(\delta).
 \label{eq:pd-closure-size}
\end{equation}
For the fixed master partition, $p_C$ is bounded below over the finitely many
rows below $S$; thus $m\ge c_{\mathcal J,S}\xi n$ for all large $n$.

We first rule out premature termination of the root-avoiding exploration.  Let
$r$ be the number of completed edges if that exploration stops because only
$v$ remains in the closure row.  Then exactly $m-1$ closure-row targets have
been used, so
\[
        R_C(r)=m-1,
        \qquad r\ge m-1.
\]
On the simultaneous prefix event of Lemma
\ref{lem:pd-two-scale-balance}, its scale-sensitive estimate at this terminal
prefix gives
\[
 \left|\frac{m-1}{r p_C}-1\right|
 \le C_{\kappa,S}\frac nr\Delta_n(U,\delta,h)
      +\frac{C_{\mathcal J}}r.
\]
Writing $\theta=r/N$ and multiplying by $\theta$ yields
\[
 \left|\frac{m-1}{Np_C}-\theta\right|
 \le C_{\kappa,S,\xi}\Delta_n(U,\delta,h)
      +\frac{C_{\mathcal J}}N.
\]
Together with \eqref{eq:pd-closure-size},
\begin{equation}
        |1-r/N|
        \le C_{\kappa,S,\xi}
        \Delta_n(U,\delta,h)+\frac{C_{\mathcal J,S,\xi}}n.
 \label{eq:pd-no-premature}
\end{equation}
If the right-hand side is smaller than $\zeta/2$, the auxiliary exploration
cannot stop before $(1-\zeta)N$ completed edges.  If it is not smaller than
$\zeta/2$, the right-hand side of \eqref{eq:pd-root-cdf}, after increasing its
constant, is at least one and the proposition is trivial.  Thus we may assume
that the auxiliary exploration is defined through every
$k\le(1-\zeta)N$.

Fix $x\in[\zeta,1-\zeta]$ and put $k=\lfloor xN\rfloor$.  The number $S(k)$
of closure opportunities is exactly $R_C(k)$ under the avoid-$v$ law.  Applying
\eqref{eq:pd-scale-sensitive-balance} at $\ell=k$, and using
$k\ge\zeta\xi n$, gives, on the simultaneous prefix event,
\[
 \left|\frac{S(k)}{k p_C}-1\right|
 \le C_{\kappa,S,\xi,\zeta}
       \Delta_n(U,\delta,h)+\frac{C_{\mathcal J,\xi,\zeta}}n.
\]
Combining this with \eqref{eq:pd-closure-size} and
$k/N=x+O(n^{-1})$ gives
\begin{equation}
 \sup_{x\in[\zeta,1-\zeta]}
 \left|\frac{S(\lfloor xN\rfloor)}m-x\right|
 \le C_{\kappa,S,\xi,\zeta}
       \Delta_n(U,\delta,h)+\frac{C_{\mathcal J,\xi,\zeta}}n.
 \label{eq:pd-closure-opportunities}
\end{equation}
The exceptional conditional probability is at most $\pi_n$.

On the event in \eqref{eq:pd-closure-opportunities}, the ratio
$S(k)/m$ lies in $[\zeta/2,1-\zeta/2]$ once the displayed error is small; in the
complementary parameter regime the final estimate is again trivial.  Equations
\eqref{eq:pd-root-rn} and \eqref{eq:pd-root-hazard} squeeze the survival
product between
\[
 \prod_{j=0}^{S(k)-1}\left(1-\frac{e^\beta}{m-j}\right)
 \quad\text{and}\quad
 \prod_{j=0}^{S(k)-1}\left(1-\frac{e^{-\beta}}{m-j}\right).
\]
Lemma \ref{lem:pd-harmonic-product}, followed by
\eqref{eq:pd-closure-opportunities}, shows that both products differ from
$1-x$ by at most
\[
 C_{\kappa,S,\xi,\zeta}
 \Delta_n(U,\delta,h)+\frac{C_{\mathcal J,\xi,\zeta}}n.
\]
On the exceptional prefix event we use the trivial bound by one.  Averaging
under $\Pp^\circ$, then adding the closure-tail probability
\eqref{eq:pd-closure-tail}, gives
\[
 \sup_{x\in[\zeta,1-\zeta]}
 \left|\Pp(L>\lfloor xN\rfloor\mid\mathcal H)-(1-x)\right|
\]
bounded by the right-hand side of \eqref{eq:pd-root-cdf}.  Since
$\{L/N\le x\}=\{L\le\lfloor xN\rfloor\}$, this is the claimed CDF estimate.
\end{proof}

\subsection{Deletion stability and regeneration}

\begin{lemma}[Deletion of an interior macroscopic cycle]
\label{lem:pd-deletion-stability}
Fix master-partition endpoints $0<V<U\le T$, $\xi>0$,
$\zeta\in(0,1/2)$, and $0\le\delta<1/4$.  On $\mathcal G_n$, start from a
$(U,\delta,\xi)$-regular history and expose the cycle of a uniform residual
root.  Let $N$ and $L$ be the residual size and the exposed cycle length.  With
$\chi_n,\pi_n$ as in Lemma \ref{lem:pd-two-scale-balance}, there is a constant
$C=C(\kappa,V,\xi,\zeta)$ such that
\begin{align}
&\Pp\Bigl(
 \zeta N\le L\le(1-\zeta)N,\ 
 \text{the post-deletion history is not }
 (V,\delta',\zeta\xi)\text{-regular}
 \ \Bigm|\ \mathcal H\Bigr)
 \le\pi_n,
 \label{eq:pd-deletion-joint}
\end{align}
where
\begin{equation}
 \delta'
 =C\bigl(\delta+\beta+\chi_n+\varepsilon_n+e^{-U/2}\bigr)
  +\frac{C_{\mathcal J,\xi,\zeta}}n.
 \label{eq:pd-deletion-error}
\end{equation}
The conclusion is asserted whenever the displayed $\delta'$ is smaller than
$1/4$, which is the only regime used below.  This is a joint-event estimate; it
does not condition on the exact random value of $L$.
\end{lemma}

\begin{proof}
Let $d_a$ be the number of vertices of the exposed cycle in $H_{n,a}$, and let
$d_{ab}$ be the number of its directed edges from $H_{n,a}$ to $H_{n,b}$.
The conditional prefix event in Lemma \ref{lem:pd-two-scale-balance} is
simultaneous over all prefix lengths.  It may therefore be evaluated at the
random closing time without conditioning on $L$.

On the event $\zeta N\le L\le(1-\zeta)N$, we have
$L\ge\zeta\xi n$.  Apply the scale-sensitive estimate
\eqref{eq:pd-scale-sensitive-balance} with inner cutoff $V$ and $\ell=L$.
For every row below $V$,
\[
 \left|\frac{d_a}{Lp_{n,a}}-1\right|
 \le C_{\kappa,V,\xi,\zeta}
       \Delta_n(U,\delta,h)+\frac{C_{\mathcal J,\xi,\zeta}}n.
\]
Regularity gives $m_a^R=Np_{n,a}(1+O(\delta))$, and hence
\begin{equation}
 \frac{d_a}{m_a^R}
 =\frac LN+O\left(
 \delta+\beta+\chi_n+\varepsilon_n+e^{-U/2}
 +\frac{C_{\mathcal J,\xi,\zeta}}n\right)
 \label{eq:pd-deleted-row-fraction}
\end{equation}
uniformly over rows below $V$.

The row-prefix estimate \eqref{eq:pd-row-prefix}, evaluated at the same random
closing time, gives
\[
 d_{ab}
 =\frac{d_a}{m_a^R}N_{ab}^R
  +O((\beta+\chi_n)N_{ab}^R).
\]
Combining this with \eqref{eq:pd-deleted-row-fraction},
\begin{equation}
 d_{ab}
 =\frac LN N_{ab}^R
  +O\left(
 \delta+\beta+\chi_n+\varepsilon_n+e^{-U/2}
 +\frac{C_{\mathcal J,\xi,\zeta}}n\right)N_{ab}^R.
 \label{eq:pd-deleted-cell-fraction}
\end{equation}
The quantity $d_{ab}$ is exactly the number of vertices removed from
$Q_{n,a}\cap H_{n,b}$, because every cycle edge from row $a$ to category $b$
uses one such target.

Write $x=L/N$ and let $R'$ be the post-deletion vertex set.  Equations
\eqref{eq:pd-deleted-row-fraction}--\eqref{eq:pd-deleted-cell-fraction} give
\[
 m_a^{R'}=(1-x)m_a^R+O(\Delta_n^*m_a^R),
 \qquad
 N_{ab}^{R'}=(1-x)N_{ab}^R+O(\Delta_n^*N_{ab}^R),
\]
where $\Delta_n^*$ denotes the parenthesized error in
\eqref{eq:pd-deleted-row-fraction}.  Since $1-x\ge\zeta$ and
$|R'|=(1-x)N\ge\zeta\xi n$, division by the new total fraction preserves the
regularity inequalities with error at most the quantity in
\eqref{eq:pd-deletion-error}.  Residual permanental factorization and
hereditary flatness remain valid after deleting the exposed edges.  The only
conditional failure event is the simultaneous prefix failure, whose
probability is at most $\pi_n$.  This proves
\eqref{eq:pd-deletion-joint}.
\end{proof}

\subsection{Independent residual fractions}

Order the cycles size-biased by choosing a uniform root from the current
residual vertex set, deleting its cycle, and repeating.  Let
\[
 V_{n,j}=\frac{\widetilde L_{n,j}}
 {n-\widetilde L_{n,1}-\cdots-\widetilde L_{n,j-1}}
\]
whenever the denominator is positive, and set $V_{n,j}=0$ after exhaustion.

\begin{lemma}[Conditional test functions from an interior CDF bound]
\label{lem:pd-cdf-test}
Let $W\in[0,1]$ and let $\mathscr F$ be a sigma-field.  Suppose that, almost
surely,
\[
        \sup_{x\in[a,b]}
        |\Pp(W\le x\mid\mathscr F)-x|\le\eta,
        \qquad 0<a<b<1.
\]
Then
\[
        \Pp(W\notin[a,b]\mid\mathscr F)
        \le a+(1-b)+2\eta.
\]
If $f$ is absolutely continuous, vanishes outside $[a,b]$, and satisfies
$f(a)=f(b)=0$, then
\[
 \left|\Ee(f(W)\mid\mathscr F)-\int_0^1f(u)\dd u\right|
 \le\eta\int_a^b|f'(u)|\dd u.
\]
\end{lemma}

\begin{proof}
The endpoint estimate follows by evaluating the CDF bound at $a$ and $b$.
For the second assertion, conditional Stieltjes integration by parts and the
vanishing boundary values give
\[
        \Ee(f(W)\mid\mathscr F)
        =-\int_a^b\Pp(W\le u\mid\mathscr F)f'(u)\dd u,
\]
whereas
\[
        \int_a^bf(u)\dd u=-\int_a^bu f'(u)\dd u.
\]
Subtracting proves the claim.
\end{proof}

\begin{proposition}[Finite-dimensional $\operatorname{GEM}(1)$ limits]
\label{prop:pd-residual-uniforms}
For every fixed $k$,
\[
        (V_{n,1},\ldots,V_{n,k})
        \Rightarrow(U_1,\ldots,U_k),
\]
where $U_1,\ldots,U_k$ are independent $\operatorname{Unif}(0,1)$ variables.
\end{proposition}

\begin{proof}
Fix $k$ and $\zeta\in(0,1/4)$, and put
$\xi_*=\zeta^{k-1}$.  We first record a noncircular parameter hierarchy.
Let $\eta>0$ be a target error.

Choose $T_0$ so large that the closure-tail term
$C_{\xi_*}e^{-T_0}$ in Proposition
\ref{prop:pd-uniform-residual-fraction} is below a prescribed small fraction of
$\eta$.  Suppose $T_0,\ldots,T_j$ have been chosen.  Let $R_j$ and $D_j$ be,
respectively, upper bounds for the principal-error coefficients in the root law
and deletion lemma with inner cutoff $T_j$, residual fraction at least
$\xi_*$, and truncation parameter $\zeta$.  These constants are now fixed.  An
error created at this level can subsequently be amplified only by
$D_{j-1},\ldots,D_0$, all of which have already been fixed.  Put
\[
        A_j=(1+R_j+D_j)\prod_{i=0}^{j-1}(1+D_i),
\]
with the empty product equal to one, and choose $T_{j+1}>T_j$ so large that
\[
        A_je^{-T_{j+1}/2}\le 2^{-j-3}\eta.
\]
Continue until $T_k$ has been chosen.

Use one master partition of $[0,T_k]$ containing
$T_0,\ldots,T_k$, append the tail interval, and let its mesh be at most $h$.
After the cutoffs are fixed, all coefficients in the finitely many root and
deletion estimates are finite.  Choose $h$ so small that the errors involving
$\beta=2\kappa h$ remain below their assigned fractions of $\eta$ after the
finite deletion recursion.  More explicitly, if $D_j$ denotes the coefficient
in \eqref{eq:pd-deletion-error} with inner cutoff $T_j$, the deterministic
recursion
\[
        d_0=0,\qquad
        d_{r+1}=D_{k-r-1}(d_r+\beta+e^{-T_{k-r}/2})
\]
has $\max_{r\le k}d_r\to0$ as all displayed input errors tend to zero.
The preceding choices make this maximum smaller than a prescribed fraction of
$\eta$, and in particular smaller than $1/8$.  Finally, for this fixed master
partition take the maxima of the finitely many sequences $\chi_n$ and $\pi_n$
that occur at the different cutoff levels, and choose $n$ so large that
\[
        \varepsilon_n+\chi_n+\pi_n+C_{\mathcal J}/n
\]
is below its assigned fraction of $\eta$ and
$\Pp(\mathcal G_n^c)\le\eta$.  This order of choices is important: the inner
cutoff is fixed before its outer cutoff, and the mesh is chosen only after all
cutoff-dependent constants are known.

On the event that
\[
        \zeta\le V_{n,j}\le1-\zeta,
        \qquad 1\le j<r,
 \tag{6.20}\label{eq:pd-interior-history}
\]
the residual size before stage $r$ is at least
$\zeta^{r-1}n\ge\xi_*n$.  At stage $r$, use outer cutoff
$T_{k-r+1}$ and inner cutoff $T_{k-r}$.  Inductively, the current history is
regular at the outer cutoff with a regularity error bounded by the deterministic
recursion above, except on a conditional event of probability $O_k(\eta)$.
Proposition \ref{prop:pd-uniform-residual-fraction} therefore gives, uniformly
over every such current history,
\begin{equation}
 \sup_{x\in[\zeta,1-\zeta]}
 \left|\Pp(V_{n,r}\le x\mid\mathscr F_{r-1})-x\right|
 \le C_k\eta.
 \label{eq:pd-stage-cdf}
\end{equation}
On the event $V_{n,r}\in[\zeta,1-\zeta]$, Lemma
\ref{lem:pd-deletion-stability} makes the next history regular at cutoff
$T_{k-r}$, except on a conditional event of probability at most $C_k\eta$.
All these estimates are uniform over the preceding regular history.

We now turn the conditional estimates into joint convergence.  Let
$f_1,\ldots,f_k$ be bounded Lipschitz functions on $[0,1]$.  Choose a piecewise
linear cutoff $\varphi_\zeta$ which is zero outside
$[\zeta,1-\zeta]$ and equal to one on
$[2\zeta,1-2\zeta]$, and put $g_j=f_j\varphi_\zeta$.  Lemma
\ref{lem:pd-cdf-test}, applied to \eqref{eq:pd-stage-cdf}, gives
\[
 \left|\Ee(g_r(V_{n,r})\mid\mathscr F_{r-1})
       -\int_0^1g_r(u)\dd u\right|
 \le C_{f_r,\zeta}\eta
\]
on every regular history.  Let $\mathcal R_r$ be the event that, after the
first $r$ deletions, the history is regular at cutoff $T_{k-r}$ with the
prescribed error; put $\mathcal R_0=\mathcal G_n$.  Because $g_r$ vanishes
outside $[\zeta,1-\zeta]$, the joint-event deletion lemma also gives, on
$\mathcal R_{r-1}$,
\[
 \left|\Ee\bigl(g_r(V_{n,r})\one_{\mathcal R_r}
       \mid\mathscr F_{r-1}\bigr)
       -\int_0^1g_r(u)\dd u\right|
 \le C_{f_r,\zeta}\eta.
\]
Multiplying by
$\prod_{j<r}g_j(V_{n,j})\one_{\mathcal R_{r-1}}$, taking expectations, and
inducting on $r$ yields
\[
 \left|\Ee\left(\prod_{j=1}^kg_j(V_{n,j})
        \one_{\mathcal R_k}\right)
       -\prod_{j=1}^k\int_0^1g_j(u)\dd u\right|
 \le C_{f,k,\zeta}\eta+o_n(1).
\]
On the event that every $g_j(V_{n,j})$ is nonzero, every residual fraction is
in the interior range, so a union bound over the conditional deletion failures
shows that removing $\one_{\mathcal R_k}$ costs another
$C_{f,k}\eta+o_n(1)$.  Consequently
\begin{equation}
 \left|\Ee\prod_{j=1}^kg_j(V_{n,j})
       -\prod_{j=1}^k\int_0^1g_j(u)\dd u\right|
 \le C_{f,k,\zeta}\eta+o_n(1).
 \label{eq:pd-truncated-product}
\end{equation}
Here the $o_n(1)$ includes $\Pp(\mathcal G_n^c)$.

Evaluating \eqref{eq:pd-stage-cdf} at $2\zeta$ and $1-2\zeta$, and adding the
regularity-failure probabilities, gives
\begin{equation}
 \limsup_{n\to\infty}\Pp\left(
 V_{n,j}\notin[2\zeta,1-2\zeta]
 \text{ for some }j\le k\right)
 \le4k\zeta+C_k\eta.
 \label{eq:pd-actual-endpoints}
\end{equation}
Consequently, replacing $g_j$ by $f_j$ in
\eqref{eq:pd-truncated-product} costs at most
$C_{f,k}(\zeta+\eta)$, and the same is true for replacing
$\int g_j$ by $\int f_j$.  Let first $n\to\infty$, then
$\eta\downarrow0$, and finally $\zeta\downarrow0$.  We obtain
\[
        \Ee\prod_{j=1}^kf_j(V_{n,j})
        \longrightarrow
        \prod_{j=1}^k\int_0^1f_j(u)\dd u.
\]
Products of bounded Lipschitz functions form a convergence-determining class on
$[0,1]^k$.  This proves convergence to $k$ independent uniform variables.
The convention $V_{n,j}=0$ after exhaustion causes no difficulty, because
exhaustion before stage $k$ is contained in the endpoint event in
\eqref{eq:pd-actual-endpoints}.
\end{proof}

\begin{proof}[Proof of Theorem \ref{thm:macroscopic-pd}]
Put
\[
 X_{n,j}=\frac{\widetilde L_{n,j}}n
 =V_{n,j}\prod_{i<j}(1-V_{n,i}),
 \qquad
 R_{n,k}=1-\sum_{j=1}^kX_{n,j}
        =\prod_{i=1}^k(1-V_{n,i}).
\]
Proposition \ref{prop:pd-residual-uniforms} and the continuous mapping theorem
give, for every fixed $k$,
\begin{align}
 &(X_{n,1},\ldots,X_{n,k},R_{n,k})\notag\\
 &\qquad\Rightarrow
 \left(U_1,(1-U_1)U_2,\ldots,
 U_k\prod_{i<k}(1-U_i),\prod_{i=1}^k(1-U_i)\right).
 \label{eq:pd-finite-gem}
\end{align}
Because $0\le R_{n,k}\le1$,
\[
        \Ee R_{n,k}\longrightarrow2^{-k}.
\]
Thus, for every $\epsilon>0$,
\[
 \limsup_{n\to\infty}\Pp(R_{n,k}>\epsilon)
 \le\epsilon^{-1}2^{-k}.
 \label{eq:pd-tail-mass}
\]
For a bounded Lipschitz functional on $\ell^1$, truncate after $k$ coordinates,
use \eqref{eq:pd-finite-gem}, and bound the truncation error by the residual
tail mass.  Equation \eqref{eq:pd-tail-mass}, followed by $k\to\infty$, proves
\[
        (X_{n,1},X_{n,2},\ldots)
        \Rightarrow\operatorname{GEM}(1)
\]
in $\ell^1$.

For finite nonnegative sequences, the usual two-point uncrossing inequality
shows that matching coordinates in decreasing order minimizes the sum of
absolute differences.  Passing to truncations gives
\[
        \|x^\downarrow-y^\downarrow\|_1\le\|x-y\|_1
\]
for nonnegative $\ell^1$ sequences.  The continuous mapping theorem therefore
gives the $\operatorname{PD}(1)$ limit.  The first residual fraction is the
cycle length of a uniform root divided by $n$, which proves the final assertion
of the theorem.

The argument was carried out along an arbitrary subsequence after extracting a
marked-profile limit.  The resulting law does not depend on that profile.
Every subsequence therefore has a further subsequence with the same limit, so
the full sequence converges.
\end{proof}

\begin{corollary}[Power sums of macroscopic cycles]
\label{cor:pd-power-sums}
Under the hypotheses of Theorem \ref{thm:macroscopic-pd}, for every $q>1$,
\[
 \sum_{j\ge1}\left(\frac{L_{n,j}}n\right)^q
 \Rightarrow\sum_{j\ge1}P_j^q,
 \qquad (P_j)\sim\operatorname{PD}(1),
\]
and
\[
 \Ee\sum_{j\ge1}\left(\frac{L_{n,j}}n\right)^q
 \longrightarrow\frac1q.
\]
\end{corollary}

\begin{proof}
The map $x\mapsto\sum_jx_j^q$ is continuous and bounded by one on
$\mathcal S^\downarrow$.  If $\widetilde P$ is a size-biased pick from a
partition $(P_j)$, then, conditionally on the partition,
\[
        \Ee(\widetilde P^{\,q-1}\mid(P_j))=\sum_jP_j^q.
\]
For a $\operatorname{PD}(1)$ partition, its size-biased ordering has the
$\operatorname{GEM}(1)$ law, so $\widetilde P$ is uniform on $(0,1)$.  Hence
\[
        \Ee\sum_jP_j^q
        =\Ee\widetilde P^{\,q-1}
        =\int_0^1u^{q-1}\dd u=1/q.
\]
Bounded convergence transfers the expectation limit.
\end{proof}

\begin{thebibliography}{9}

\bibitem{Pitman2006}
J. Pitman.
\newblock \emph{Combinatorial Stochastic Processes}.
\newblock Lecture Notes in Mathematics, vol. 1875.
\newblock Springer-Verlag, Berlin, 2006.

\end{thebibliography}

\end{document}
